\textbf{Answer: 1)} $v = C\sqrt{GM_{\rm E}/R}$

\underline{\textbf{Solution for 1}}


It is tempting to use the same reasoning as in the low-orbit satellite example in Section 1.2. This reasoning gives the same result, $v = C\sqrt{gR} = C\sqrt{GM_{\rm E}/R}$, where $C$ is some number (it turns out that $C = \sqrt{2}$). Although this solution yields the correct answer, it isn't quite rigorous, in view of the footnote in the low-orbit satellite example. Because the particle isn't always at the same radius, the force changes, so it isn't obvious that we can absorb the $M_{\rm E}$ and $G$ dependence into one quantity, $g$, as we did with the orbiting satellite. Let us therefore be more rigorous with the following reasoning.

The dimensionful quantities in the problem are $[m] = M$, the radius of the earth $[R] = L$, the mass of the earth $[M_{\rm E}] = M$, and Newton's gravitational constant $[G] = L^3/MT^2$. These units for $G$ follow from the gravitational force law, $F = Gm_1m_2/r^2$. If we use no information other than these given quantities, then there is no way to arrive at the speed of $C\sqrt{GM_{\rm E}/R}$, because for all we know, there could be a factor of $(m/M_{\rm E})^7$ in the answer. This number is dimensionless, so it wouldn't mess up the units.

If we want to make any progress in this problem, we have to use the fact that the gravitational force takes the form of $GM_{\rm E}m/r^2$. This then implies (as was stated in the problem) that the acceleration is independent of $m$. And since the path of the particle is determined by its acceleration, we see that the answer can't depend on $m$. We are therefore left with the quantities $G$, $R$, and $M_{\rm E}$, and you can show that the only combination of these quantities that gives the units of speed is $v = C\sqrt{GM_{\rm E}/R}$.